% 20-07(20쪽) — 구간 [-3,2] 에서 정의된 y=f(x) 의 그래프. x<-2 는 y=2, -2<x<0 은 기울기 -1/2 인 직선, x>0 은 기울기 -2 인 직선.
% 원문의 빈 점·찬 점과 점선(-3·-2 세로선, 1·2·-2 가로선)만 옮긴다. 스캔 화질이 낮아 TikZ 로 다시 그림.
\begin{tikzpicture}[line width=0.55pt, x=0.58cm, y=0.58cm, every node/.style={font=\footnotesize}]
  \draw[->, >=stealth, line width=0.7pt] (-3.9,0) -- (3.1,0) node[below=1pt] {$x$};
  \draw[->, >=stealth, line width=0.7pt] (0,-3.4) -- (0,2.9) node[left=1pt] {$y$};
  % 점선
  \draw[densely dashed, line width=0.4pt] (-3,0) -- (-3,2) (-2,0) -- (-2,1);
  \draw[densely dashed, line width=0.4pt] (-2,2) -- (0,2) (-2,1) -- (0,1);
  \draw[densely dashed, line width=0.4pt] (0,-2) -- (2,-2) (2,0) -- (2,-2);
  % x<-2 : y=2
  \draw[line width=0.6pt] (-3.8,2) -- (-2,2);
  % -2<x<0 : 기울기 -1/2
  \draw[line width=0.6pt] (-2,2) -- (0,1);
  % x>0 : 기울기 -2
  \draw[line width=0.6pt] (0,2) -- (2.5,-3);
  % 빈 점
  \draw[fill=white, line width=0.5pt] (-2,2) circle (2.2pt);
  \draw[fill=white, line width=0.5pt] (0,1) circle (2.2pt);
  % 찬 점
  \fill (-2,1) circle (2.2pt);
  \fill (0,2) circle (2.2pt);
  % 라벨
  \node[anchor=west, inner sep=2.5pt] at (0.14,2) {$2$};
  \node[anchor=north east, inner sep=1.5pt] at (-0.1,0.9) {$1$};
  \node[anchor=east, inner sep=2.5pt, fill=white] at (-0.1,-2) {$-2$};
  \node[below=1pt] at (-3,-0.05) {$-3$};
  \node[below=1pt] at (-2,-0.05) {$-2$};
  \node[below right=-2pt and -3pt] at (0,0) {$\pt{O}$};
  \node[below=1pt] at (1,-0.05) {$1$};
  \node[above=2pt] at (2,0.05) {$2$};
  \node[anchor=west] at (1.05,1.3) {$y=f(x)$};
\end{tikzpicture}
